%% Phase-4 migration of v1.1 Part III, source lines 3022--3926.
%% Body statements only; every complete source proof is compressed in
%% annex/p02/first-order-syntax-semantics.tex.
%% R-010 is preserved in the total substitution definition and its two
%% supporting lemmas.  R-024 is preserved in the definitional-extension entry.

\chapter{First-Order Syntax and Semantics}
\label{sec:p02-first-order-syntax-semantics}

First-order logic adds variables over a nonempty domain, function and relation
symbols, and quantifiers.  Equality below is logical identity, not an
unconstrained binary relation.  The conventions follow the standard set-based
accounts \cite{EbbinghausFlumThomas1994,Hodges1993,ChiswellHodges2007}; the
handbook notation replaces structure-domain bars by \(\dom\mathcal M\).

\section{Signatures, terms, and formulas}
\label{sec:p02-fol-grammar}

\begin{hbdefinition}{def:p02-fol-signature}{Signature, first-order@First-order signature}
Fix the countably infinite object-variable set
\(\mathsf{Var}\defeq\{x_n:n\in\N\}\).  Letters \(x,y,z\) range over it;
\(r,t,u\) range over terms; and \(s\) is reserved for assignments.  A
\emph{first-order signature}\index{signature!first-order} is
\[
 \Lang=\langle\mathcal F,\mathcal R,
       \operatorname{ar}_{\mathcal F},\operatorname{ar}_{\mathcal R}\rangle,
\]
where \(\operatorname{ar}_{\mathcal F}:\mathcal F\to\N\) and
\(\operatorname{ar}_{\mathcal R}:\mathcal R\to\N\).  A zero-ary function
symbol is a constant.  Equality is logical and does not belong to
\(\mathcal R\).  The sets \(\mathcal F,\mathcal R,\mathsf{Var}\), and
\(\{=,\bottom,\neg,\land,\lor,\to,\forall,\exists,
\texttt{(},\texttt{)},\texttt{,}\}\) are pairwise disjoint.
The signature cardinality is
\(\abs{\Lang}\defeq\abs{\mathcal F\cup\mathcal R}\): arity data introduce no
symbols.  A signature is syntactic; a structure later supplies
interpretations.
\end{hbdefinition}

\begin{hbdefinition}{def:p02-fol-term}{Term, first-order@First-order term}
For \(\Lang\), \(\operatorname{Term}(\Lang)\) is the least set containing
every \(x\in\mathsf{Var}\) and containing
\(f(t_1,\ldots,t_n)\) whenever \(f\in\mathcal F\) is \(n\)-ary and all
\(t_i\) are terms.  When \(n=0\), write \(f\), not \(f()\).  No other
string is a term.\index{term!first-order}
\end{hbdefinition}

\begin{hbexample}{def:p02-fol-arithmetic-terms}{Arithmetic terms}
For function symbols \(0,S,+,\cdot\) of arities \(0,1,2,2\), respectively,
and binary relation symbol \(<\), the following are terms:
\[
 x_0,\qquad S(0),\qquad +(S(0),x_2),\qquad
 \cdot(+(x_0,x_1),S(x_2)).
\]
The string \(<(x_0,x_1)\) is not a term.  Infix notation such as
\(x_0+S(0)\) abbreviates the prefix grammar.
\end{hbexample}

\begin{hbdefinition}{def:p02-fol-formula}{Formula, first-order@First-order formula}
The set \(\operatorname{Fm}(\Lang)\) is the least set generated by
\[
\begin{array}{ll}
 r=t,&R(t_1,\ldots,t_n),\\
 \bottom,&\neg\varphi,\\
 (\varphi\circ\psi)&(\circ\in\{\land,\lor,\to\}),\\
 \forall x\,\varphi,&\exists x\,\varphi,
\end{array}
\]
where all displayed terms are in \(\operatorname{Term}(\Lang)\),
\(R\in\mathcal R\) has the displayed arity, and
\(x\in\mathsf{Var}\).  The first two forms are \emph{atomic}; a zero-ary
relation symbol is written \(R\), not \(R()\).  No other string is a
formula.\index{formula!first-order}\index{atomic formula}
\end{hbdefinition}

\begin{hbconvention}{def:p02-fol-quantifier-scope}{Quantifier scope}
In the official grammar a quantifier governs the formula immediately after
its variable.  In abbreviated displays it binds more tightly than every
binary connective, and parentheses extend its scope.  Thus
\(\forall x\,P(x)\to Q(x)\) abbreviates
\((\forall x\,P(x))\to Q(x)\), not
\(\forall x(P(x)\to Q(x))\).
\end{hbconvention}

\begin{hbtheorem}{thm:p02-fol-unique-readability}{proved=annex:p2.2.syn.1}{Unique readability for first-order expressions}
Every \(\Lang\)-term has exactly one outermost formation clause and uniquely
determined immediate subterms.  Every \(\Lang\)-formula has exactly one
outermost clause and uniquely determined principal symbol, bound variable
when present, and immediate subformulas.
\end{hbtheorem}

\section{Variables, complexity, and substitution}
\label{sec:p02-fol-substitution}

\begin{hbdefinition}{def:p02-fol-occurring-variables}{Variable occurrence sets}
The sans-serif \(\mathsf{Var}\) is the fixed variable alphabet; the operator
\(\operatorname{Var}\) records occurrences.  For terms,
\[
 \operatorname{Var}(x)=\{x\},\qquad
 \operatorname{Var}(f(t_1,\ldots,t_n))
   =\bigcup_{i=1}^{n}\operatorname{Var}(t_i).
\]
For formulas, \(\operatorname{Var}(\varphi)\) contains every variable with a
free or bound occurrence, recursively:
\[
\begin{aligned}
 \operatorname{Var}(r=t)&=\operatorname{Var}(r)\cup\operatorname{Var}(t),&
 \operatorname{Var}(R(t_1,\ldots,t_n))&=\bigcup_{i=1}^n\operatorname{Var}(t_i),\\
 \operatorname{Var}(\bottom)&=\varnothing,&
 \operatorname{Var}(\neg\varphi)&=\operatorname{Var}(\varphi),\\
 \operatorname{Var}(\varphi\circ\psi)&=\operatorname{Var}(\varphi)\cup
 \operatorname{Var}(\psi),&
 \operatorname{Var}(Qx\,\varphi)&=\operatorname{Var}(\varphi)\cup\{x\}.
\end{aligned}
\]
\index{variable!occurrence set}
\end{hbdefinition}

\begin{hbdefinition}{def:p02-fol-free-variables}{Free variables and sentences}
The free-variable operation is
\[
\begin{aligned}
 \FV(r=t)&=\operatorname{Var}(r)\cup\operatorname{Var}(t),&
 \FV(R(t_1,\ldots,t_n))&=\bigcup_{i=1}^n\operatorname{Var}(t_i),\\
 \FV(\bottom)&=\varnothing,&
 \FV(\neg\varphi)&=\FV(\varphi),\\
 \FV(\varphi\circ\psi)&=\FV(\varphi)\cup\FV(\psi),&
 \FV(Qx\,\varphi)&=\FV(\varphi)\setminus\{x\},
\end{aligned}
\]
for \(\circ\in\{\land,\lor,\to\}\) and
\(Q\in\{\forall,\exists\}\).  A formula with no free variable is a
\emph{sentence}; the set of sentences is
\(\operatorname{Sent}(\Lang)\).  For a set \(\Gamma\) of formulas,
\(\FV(\Gamma)\defeq\bigcup_{\gamma\in\Gamma}\FV(\gamma)\).
\index{free variable}\index{sentence!first-order}
\end{hbdefinition}

\begin{hbexample}{def:p02-fol-free-bound-occurrence}{Free and bound occurrences}
An occurrence of \(x\) is bound precisely when it lies in the scope of a
governing \(\forall x\) or \(\exists x\).  In
\(P(x)\land\exists x\,Q(x,y)\), the first \(x\) is free, the second is
bound, and \(y\) is free.  Quantifiers bind occurrences, not every instance
of the same variable token.
\end{hbexample}

\begin{hbdefinition}{def:p02-fol-complexity}{Complexity, first-order formula@First-order formula complexity}
Set
\[
\begin{aligned}
 c(r=t)&=c(R(t_1,\ldots,t_n))=c(\bottom)=0,\\
 c(\neg\varphi)&=c(\varphi)+1,\\
 c(\varphi\circ\psi)&=c(\varphi)+c(\psi)+1,\\
 c(Qx\,\varphi)&=c(\varphi)+1.
\end{aligned}
\]
Every immediate subformula has smaller complexity.  All recursion or
induction below described as being on formula complexity uses this measure.
\index{formula complexity!first-order}
\end{hbdefinition}

\begin{hbdefinition}{def:p02-fol-term-substitution}{Substitution in a term}
For terms \(r,t\) and variable \(x\), define
\[
\begin{aligned}
 x[t/x]&=t,&y[t/x]&=y\quad(y\ne x),\\
 f(r_1,\ldots,r_n)[t/x]
   &=f(r_1[t/x],\ldots,r_n[t/x]).
\end{aligned}
\]
Only occurrences of \(x\) are replaced.\index{substitution!in a term}
\end{hbdefinition}

\begin{hbdefinition}{def:p02-fol-capture-avoiding-substitution}{Substitution, capture-avoiding@Capture-avoiding substitution}
Order the variables as \(x_0,x_1,\ldots\).  For a formula \(\varphi\), term
\(t\), and variable \(x\), define \(\varphi[t/x]\) by recursion on formula
complexity: substitute in every atomic term and commute with
\(\bottom,\neg,\land,\lor,\to\).  For
\(Q\in\{\forall,\exists\}\), put
\[
 (Qy\,\psi)[t/x]=
 \begin{cases}
  Qy\,\psi,&y=x,\\
  Qy(\psi[t/x]),&y\ne x\text{ and }y\notin\operatorname{Var}(t),\\
  Qz((\psi[z/y])[t/x]),&y\ne x\text{ and }y\in\operatorname{Var}(t),
 \end{cases}
\]
where in the last case \(z\) is the least variable outside
\(\operatorname{Var}(\psi)\cup\operatorname{Var}(t)\cup\{x,y\}\).
This is a simultaneous strong recursion: variable-for-variable substitution
preserves complexity, so
\(c(\psi[z/y])=c(\psi)<c(Qy\,\psi)\).  The final clause first renames the
binder and then substitutes; it prevents capture of a free variable of
\(t\).\index{substitution!capture-avoiding}
\end{hbdefinition}

\begin{hblemma}{thm:p02-fol-substitution-complexity}{proved=annex:p2.2.syn.2}{Substitution preserves formula complexity}
For every formula \(\varphi\), term \(t\), and variable \(x\),
\(c(\varphi[t/x])=c(\varphi)\).
\end{hblemma}

\begin{hbdefinition}{def:p02-fol-alpha-equivalence}{Alpha-equivalence}
Alpha-equivalence, written \(=_{\alpha}\), is the least equivalence relation
on formulas that is a congruence for every formula constructor and contains
\[
 Qx\,\varphi =_{\alpha} Qy(\varphi[y/x])
\]
for \(Q\in\{\forall,\exists\}\) whenever
\(y\notin\operatorname{Var}(\varphi)\).  A generating step renames one bound
variable and exactly the occurrences governed by its quantifier.  It is
syntactic, not literal string identity and not logical equivalence.
\index{alpha-equivalence}
\end{hbdefinition}

\begin{hbexample}{def:p02-fol-avoiding-capture}{Avoiding capture}
If \(\varphi=\forall y\,R(x,y)\), textual replacement of \(x\) by \(y\)
would incorrectly give \(\forall y\,R(y,y)\).  If \(z\) is the least allowed
fresh variable, the defined result is
\(\varphi[y/x]=\forall z\,R(y,z)\), and the inserted \(y\) stays free.
\end{hbexample}

\begin{hblemma}{thm:p02-fol-free-variables-substitution}{proved=annex:p2.2.syn.3}{Free variables under substitution}
For terms \(r,t\) and variable \(x\), if
\(x\notin\operatorname{Var}(r)\), then \(r[t/x]=r\), and always
\[
 \operatorname{Var}(r[t/x])\subseteq
 (\operatorname{Var}(r)\setminus\{x\})\cup\operatorname{Var}(t).
\]
For every formula \(\varphi\), term \(t\), and variables \(x,y\):
\begin{enumerate}
 \item \(\FV(\varphi[t/x])\subseteq
 (\FV(\varphi)\setminus\{x\})\cup\operatorname{Var}(t)\);
 \item if \(y\notin\operatorname{Var}(\varphi)\), then
 \[
 \FV(\varphi[y/x])=
 \begin{cases}
  (\FV(\varphi)\setminus\{x\})\cup\{y\},&x\in\FV(\varphi),\\
  \FV(\varphi),&x\notin\FV(\varphi).
 \end{cases}
 \]
\end{enumerate}
Term variables are measured by \(\operatorname{Var}\); no free-variable
operator on terms is being introduced.
\end{hblemma}

\begin{hblemma}{thm:p02-fol-inert-substitution}{proved=annex:p2.2.syn.4}{Substitution for a nonfree variable}
If \(x\notin\FV(\varphi)\), then
\(\varphi[t/x]=_{\alpha}\varphi\).
\end{hblemma}

The first-order subformula-instance relation and its exact cut-free
subformula property are recorded in
Section~\ref{sec:p02-fol-subformula-property}.  Literal propositional
subformulahood alone is insufficient because quantified rules introduce term
instances and fresh eigenvariables.

\section{Structures, assignments, and satisfaction}
\label{sec:p02-fol-semantics}

\begin{hbdefinition}{def:p02-fol-structure}{Structure, first-order@First-order structure}
A \(\Lang\)-structure \(\mathcal M\) consists of a nonempty domain
\(\dom\mathcal M\), a total function
\(f^{\mathcal M}:(\dom\mathcal M)^n\to\dom\mathcal M\) for each \(n\)-ary
\(f\in\mathcal F\), and a relation
\(R^{\mathcal M}\subseteq(\dom\mathcal M)^n\) for each \(n\)-ary
\(R\in\mathcal R\).  A constant \(c\) is interpreted by
\(c^{\mathcal M}\in\dom\mathcal M\).  Equality is actual identity on
\(\dom\mathcal M\).\index{structure!first-order}\index{domain!of a structure}
\end{hbdefinition}

\begin{hbdefinition}{def:p02-fol-expansion-reduct}{Signature extension, expansion, and reduct}
A signature \(\Lang'\) \emph{extends} \(\Lang\) when its function and
relation-symbol sets include those of \(\Lang\) and its arity maps restrict
to the old maps.  A \(\Lang'\)-structure \(\mathcal M'\) is an
\emph{expansion} of a \(\Lang\)-structure \(\mathcal M\) when the domains
and all old-symbol interpretations agree.  The \emph{reduct}
\(\mathcal M'\mathbin{\upharpoonright}\Lang\) forgets the new
interpretations.  Neither operation changes the domain, so neither means
substructure.\index{expansion}\index{reduct}
\end{hbdefinition}

\begin{hbdefinition}{def:p02-fol-assignment}{Assignment and update}
An assignment in \(\mathcal M\) is a function
\(s:\mathsf{Var}\to\dom\mathcal M\).  For \(a\in\dom\mathcal M\),
\[
 s[x\mapsto a](y)=
 \begin{cases}a,&y=x,\\s(y),&y\ne x.\end{cases}
\]
\index{assignment!first-order}\index{assignment!update}
\end{hbdefinition}

\begin{hbdefinition}{def:p02-fol-term-value}{Term value}
The value of a term under \(s\) is defined by
\[
 \operatorname{val}^{\mathcal M}_s(x)=s(x),\qquad
 \operatorname{val}^{\mathcal M}_s(f(t_1,\ldots,t_n))
 =f^{\mathcal M}(\operatorname{val}^{\mathcal M}_s(t_1),\ldots,
                  \operatorname{val}^{\mathcal M}_s(t_n)).
\]
\index{term!value}
\end{hbdefinition}

\begin{hblemma}{thm:p02-fol-term-coincidence}{proved=annex:p2.2.syn.5}{Term coincidence}
If assignments \(s,s'\) agree on \(\operatorname{Var}(t)\), then
\(\operatorname{val}^{\mathcal M}_s(t)
=\operatorname{val}^{\mathcal M}_{s'}(t)\).
\end{hblemma}

\begin{hblemma}{thm:p02-fol-term-substitution}{proved=annex:p2.2.syn.6}{Term-substitution lemma}
For all terms \(r,t\), assignments \(s\), and variables \(x\),
\[
\begin{aligned}
 \operatorname{val}^{\mathcal M}_s(r[t/x])
 &=\operatorname{val}^{\mathcal M}_{s'}(r),\\
 s'&=s[x\mapsto\operatorname{val}^{\mathcal M}_s(t)].
\end{aligned}
\]
\end{hblemma}

\begin{hbdefinition}{def:p02-fol-satisfaction}{Satisfaction, first-order@First-order satisfaction}
For a structure \(\mathcal M\), assignment \(s\), and formula \(\varphi\),
define \(\mathcal M,s\entails\varphi\) recursively by
\[
\begin{aligned}
 \mathcal M,s\entails r=t
 &\iff\operatorname{val}^{\mathcal M}_s(r)
       =\operatorname{val}^{\mathcal M}_s(t),\\
 \mathcal M,s\entails R(t_1,\ldots,t_n)
 &\iff(\operatorname{val}^{\mathcal M}_s(t_1),\ldots,
        \operatorname{val}^{\mathcal M}_s(t_n))\in R^{\mathcal M},\\
 \mathcal M,s&\not\entails\bottom,\\
 \mathcal M,s\entails\neg\varphi
 &\iff\mathcal M,s\not\entails\varphi,\\
 \mathcal M,s\entails\varphi\land\psi
 &\iff\mathcal M,s\entails\varphi\text{ and }\mathcal M,s\entails\psi,\\
 \mathcal M,s\entails\varphi\lor\psi
 &\iff\mathcal M,s\entails\varphi\text{ or }\mathcal M,s\entails\psi,\\
 \mathcal M,s\entails\varphi\to\psi
 &\iff\mathcal M,s\not\entails\varphi\text{ or }\mathcal M,s\entails\psi,\\
 \mathcal M,s\entails\forall x\,\varphi
 &\iff\mathcal M,s[x\mapsto a]\entails\varphi
       \text{ for every }a\in\dom\mathcal M,\\
 \mathcal M,s\entails\exists x\,\varphi
 &\iff\mathcal M,s[x\mapsto a]\entails\varphi
       \text{ for some }a\in\dom\mathcal M.
\end{aligned}
\]
\index{satisfaction!first-order}
\end{hbdefinition}

\begin{hblemma}{thm:p02-fol-coincidence}{proved=annex:p2.2.syn.7}{Formula coincidence}
If assignments \(s,s'\) agree on \(\FV(\varphi)\), then
\(\mathcal M,s\entails\varphi\) iff
\(\mathcal M,s'\entails\varphi\).
\end{hblemma}

\begin{hbtheorem}{thm:p02-fol-substitution}{proved=annex:p2.2.syn.8}{Substitution lemma}
For every formula \(\varphi\), term \(t\), variable \(x\), structure
\(\mathcal M\), and assignment \(s\),
\[
 \mathcal M,s\entails\varphi[t/x]
 \quad\Longleftrightarrow\quad
 \mathcal M,s[x\mapsto\operatorname{val}^{\mathcal M}_s(t)]
 \entails\varphi.
\]
\end{hbtheorem}

\begin{hbcorollary}{thm:p02-fol-alpha-invariance}{proved=annex:p2.2.syn.9}{Alpha-invariance}
If \(\varphi=_{\alpha}\psi\), then
\(\mathcal M,s\entails\varphi\) iff
\(\mathcal M,s\entails\psi\) for every structure \(\mathcal M\) and
assignment \(s\).
\end{hbcorollary}

\begin{hbcorollary}{thm:p02-fol-sentence-independence}{proved=annex:p2.2.syn.10}{Assignment independence for sentences}
If \(\sigma\in\operatorname{Sent}(\Lang)\), then
\(\mathcal M,s\entails\sigma\) iff
\(\mathcal M,s'\entails\sigma\) for all assignments \(s,s'\) in
\(\mathcal M\).
\end{hbcorollary}

\begin{hbdefinition}{def:p02-fol-semantic-notions}{Truth, validity, consequence, and formula equivalence}
For a sentence \(\sigma\), write \(\mathcal M\entails\sigma\) when
\(\mathcal M,s\entails\sigma\) for one, equivalently every, assignment.
For \(\Gamma\subseteq\operatorname{Fm}(\Lang)\),
\(\Gamma\entails\varphi\) means that for every \(\Lang\)-structure
\(\mathcal M\) and assignment \(s\), satisfaction of every
\(\gamma\in\Gamma\) implies satisfaction of \(\varphi\).  A formula is
\emph{valid} when \(\varnothing\entails\varphi\).  Formula equivalence is
\[
 \varphi\Leftrightarrow\psi
 \quad\Longleftrightarrow\quad
 \text{for all \(\mathcal M,s\), }
 \mathcal M,s\entails\varphi\text{ iff }\mathcal M,s\entails\psi.
\]
It is neither alpha-equivalence nor literal identity.
\index{truth!in a structure}\index{validity!first-order}
\index{semantic consequence!first-order}\index{formula equivalence}
\end{hbdefinition}

\section{Translation with quantifiers}
\label{sec:p02-fol-translation}

\begin{hbconvention}{def:p02-fol-restricted-quantifiers}{Restricted quantification}
Every translation declares its signature and variable range.  If unary
\(A\) expresses the intended restriction, then
\[
\begin{aligned}
 \forall x\in A\,\varphi&\quad\text{abbreviates}\quad
 \forall x(A(x)\to\varphi),\\
 \exists x\in A\,\varphi&\quad\text{abbreviates}\quad
 \exists x(A(x)\land\varphi).
\end{aligned}
\]
\index{quantifier!restricted}
\end{hbconvention}

\begin{hbexample}{def:p02-fol-quantifier-order}{Quantifier order and witness dependence}
With \(S,P,A\) meaning student, professor, and admires,
\[
 \forall x(S(x)\to\exists y(P(y)\land A(x,y)))
\]
allows a professor to depend on the student, whereas
\[
 \exists y(P(y)\land\forall x(S(x)\to A(x,y)))
\]
requires one professor for all students.  The second implies the first, but
not conversely.
\end{hbexample}

\begin{hbexample}{def:p02-fol-negated-universal}{Negated universal and vacuity}
``No professor admires every student'' is
\(\forall y(P(y)\to\neg\forall x(S(x)\to A(y,x)))\), equivalently in
classical logic
\(\forall y(P(y)\to\exists x(S(x)\land\neg A(y,x)))\).  With no students,
both formulations require that there be no professors: the inner universal
is vacuously true and the existential false.
\end{hbexample}

\section{Theories and axiomatization}
\label{sec:p02-fol-theories}

\begin{hbdefinition}{def:p02-fol-theory}{Theory, model, satisfiability, and structure theory}
A \emph{theory} in \(\Lang\) is a set
\(T\subseteq\operatorname{Sent}(\Lang)\).  A structure \(\mathcal M\) is a
\emph{model} of \(T\), written \(\mathcal M\entails T\), iff it satisfies
every member; the model class is \(\operatorname{Mod}(T)\).  A set
\(\Gamma\subseteq\operatorname{Fm}(\Lang)\) is \emph{satisfiable} iff some
\(\mathcal M,s\) satisfies every member.  Sentence independence makes a
theory satisfiable exactly when it has a model.  The structure theory is
\[
 \operatorname{Th}_{\Lang}(\mathcal M)
 \defeq\{\sigma\in\operatorname{Sent}(\Lang):
             \mathcal M\entails\sigma\}.
\]
When \(\Lang\) is fixed, \(\operatorname{Th}(\mathcal M)\) abbreviates it.
An arbitrary theory need not be semantically closed; its closure is
\[
 \operatorname{Cn}(T)\defeq
 \{\sigma\in\operatorname{Sent}(\Lang):T\entails\sigma\}.
\]
\index{theory!first-order}\index{model!of a theory}
\index{satisfiable!first-order}\index{theory!of a structure}
\index{semantic closure}
\end{hbdefinition}

\begin{hbdefinition}{def:p02-fol-axiomatization}{Axiomatization}
A theory \(T\) \emph{axiomatizes} a class \(\mathcal K\) of
\(\Lang\)-structures iff \(\operatorname{Mod}(T)=\mathcal K\).  The class
is \emph{finitely axiomatizable} iff some finite theory axiomatizes it.
\index{axiomatization}\index{finitely axiomatizable}
\end{hbdefinition}

For \(n\ge2\), put
\[
 \iota_n\defeq\exists x_0\cdots\exists x_{n-1}
 \bigwedge_{0\le i<j<n}x_i\ne x_j,
\]
with pairs ordered lexicographically under the global finite-conjunction
convention; \(r\ne t\) abbreviates \(\neg(r=t)\).

\begin{table}[tp]
\centering
\caption{Axiomatization patterns from the verified source.}
\label{tab:p02-fol-axiomatizations}
\begin{tabular}{p{.20\textwidth}p{.70\textwidth}}
\toprule
Class & Axioms or schema \\
\midrule
Groups & Associativity; two-sided identity \(e\); two-sided inverse
\(i(x)\), in the signature \(\{\cdot,i,e\}\). Equality remains identity.\\
Dense linear orders without endpoints & Irreflexivity, transitivity,
linearity, density, and existence of a point on either side of every point,
in the signature \(\{<\}\).\\
Infinite structures & The theory \(\{\iota_n:n\ge2\}\) axiomatizes the
infinite structures.  This is a set of sentences, not one first-order
sentence.\\
\bottomrule
\end{tabular}
\end{table}

%% Depth row D-P3-010 (E): full definition, theorem, and annex proof.
\begin{hbdefinition}{def:p02-fol-definitional-extension}{Definitional extension}
Let \(T\) be a \(\Lang\)-theory; let \(R\notin\Lang\) be an \(n\)-ary
relation symbol; let \(x_1,\ldots,x_n\) be pairwise distinct; and let
\(\varphi(x_1,\ldots,x_n)\) be a \(\Lang\)-formula satisfying
\(\FV(\varphi)\subseteq\{x_1,\ldots,x_n\}\).  Let \(\Lang_R\) be the
signature obtained by adjoining \(R\) to \(\Lang\) with arity \(n\), and put
\[
 \delta_R\defeq\forall x_1\cdots\forall x_n
 \bigl(R(x_1,\ldots,x_n)\leftrightarrow
       \varphi(x_1,\ldots,x_n)\bigr).
\]
When \(n=0\), the atomic formula in this schema is the nullary symbol \(R\),
in accordance with the formula grammar above.
The passage from \(T\) to \(T\cup\{\delta_R\}\) is a
\emph{definitional extension}: it names the old-signature relation defined by
\(\varphi\), rather than adding a constraint on old symbols.
\index{definitional extension}\index{defining axiom}
\end{hbdefinition}

\begin{hbtheorem}{thm:p02-definitional-conservativity}{proved=annex:p2.2.syn.11}{Conservativity of definitional extensions}
Under Definition~\ref{def:p02-fol-definitional-extension}:
\begin{enumerate}
 \item every \(\Lang\)-model \(\mathcal M\) of \(T\) has exactly one
 \(\Lang_R\)-expansion satisfying \(\delta_R\), and it is a model of
 \(T\cup\{\delta_R\}\);
 \item for every \(\sigma\in\operatorname{Sent}(\Lang)\),
 \[
  T\cup\{\delta_R\}\entails\sigma
  \quad\Longleftrightarrow\quad T\entails\sigma.
 \]
\end{enumerate}
Consequently, once soundness and the stated Axiom-of-Choice form of
first-order completeness are available, the same equivalence holds with
\(\entails\) replaced by
\(\derives_{\mathsf{ND}_{\mathsf{FOL}}}\).  The reverse derivability
direction also follows directly by weakening the assumption set.
\end{hbtheorem}

\begin{hbwarning}{def:p02-fol-axiomatization-errors}{Axiomatization checklist}
Check that every displayed symbol belongs to the declared signature, every
arity is respected, every axiom is a sentence, and the nonempty-domain
convention is accounted for.  Universal restrictions use implications and
existential restrictions use conjunctions.  Identify a schema as a set of
sentences, never as one first-order sentence.  Finally establish both
directions of \(\operatorname{Mod}(T)=\mathcal K\).
\end{hbwarning}
